\chapter{Results}
\label{ch:results}

% Replace this paragraph with a summary of your findings. Cite
% \cite{openlab2024} where appropriate.
The evaluation was run on a small public sample data set~\cite{openlab2024}
and produced the values summarised in this chapter. Every number reported
here comes from the same run, so the procedure can be repeated
independently.

% Graphviz measurement workflow diagram. Replace the graph with your own
% workflow. Keep style= and the caption (wrap a caption containing a comma
% in braces).
\begin{graphviz}[style=editorial, width=0.8\linewidth, pos=H, caption={Measurement workflow}]
  digraph G {
    rankdir=LR
    measure -> validate
    validate -> report
  }
\end{graphviz}

% Replace this paragraph with a sentence linking the diagram to the code
% example. Diagrams carry no label, so refer to them by name rather than
% with \ref.
The workflow above measures each case, validates the measurement and
reports it; the listing below performs the normalisation step of that
workflow.

% Python code block (at most fifteen lines, numbers=true). Replace it with
% your own snippet. A listing is not a figure or a table, so it takes no
% label and appears in neither list.
\begin{code}[lang=python, numbers=true]
def normalize(values):
    """Return the values scaled to the unit interval."""
    if not values:
        return []
    low, high = min(values), max(values)
    if high == low:
        return [0.0 for _ in values]
    return [(v - low) / (high - low) for v in values]
\end{code}

% Replace with a short sentence describing what the code does.
The function rescales the measured values so that the smallest becomes zero
and the largest becomes one, which makes the cases comparable.
